\documentclass[11pt,a4paper]{article}

% --- Basics ---
\usepackage[utf8]{inputenc}
\usepackage[T1]{fontenc}
\usepackage[ngerman]{babel}
\usepackage{lmodern}
\usepackage{microtype}
\usepackage{geometry}
\geometry{margin=2.6cm}

% --- Math ---
\usepackage{amsmath,amssymb,amsthm,mathtools}
\usepackage{bm}

% --- Links ---
\usepackage{hyperref}

% --- Theorem environments (deutsch) ---
\theoremstyle{plain}
\newtheorem{satz}{Satz}
\newtheorem{lemma}{Lemma}
\newtheorem{korollar}{Korollar}

\theoremstyle{definition}
\newtheorem{definition}{Definition}
\newtheorem{axiom}{Axiom}
\newtheorem{conjecture}{Vermutung}

\theoremstyle{remark}
\newtheorem{bemerkung}{Bemerkung}

% --- Convenience ---
\newcommand{\C}{\mathbb{C}}
\newcommand{\R}{\mathbb{R}}
\newcommand{\Hh}{\mathbb{H}}
\newcommand{\Okt}{\mathbb{O}}
\newcommand{\ii}{\mathrm{i}}
\newcommand{\dd}{\mathrm{d}}
\newcommand{\norm}[1]{\left\lVert #1 \right\rVert}

\title{%
Die kritische Achse als Dualit\"ats-Fixpunkt\\
in den normierten Divisionsalgebren $\R,\C,\Hh,\Okt$\\
(Paper A: mathematischer Satz + klar markierte Vermutung zur RH)
}
\author{Thomas Hoffbauer\\Bamberg, Deutschland}
\date{12. Februar 2026}

\begin{document}
\maketitle

\begin{abstract}
Wir isolieren einen rein mathematischen Kern: In jeder normierten Divisionsalgebra
$\mathbb{A}\in\{\R,\C,\Hh,\Okt\}$ erzwingt die Norm-Dualit\"atsbedingung
$\|x\|=\|1-x\|$ f\"ur ein Element $x$ in einem komplexen Slice $\C_u\subset\mathbb{A}$
notwendig die Fixpunktlage $\Re(x)=1/2$.
Dieser Satz ist elementar, aber strukturell bedeutsam: Er identifiziert die ``Einhalb'' als
Symmetrieachse einer unitarit\"atsvertr\"aglichen Dualit\"at, unabh\"angig von Physik.

Die Verbindung zur Riemannschen Vermutung wird strikt getrennt: Wir formulieren eine
Vermutung, dass die nichttrivialen Nullstellen der (vollendeten) Zetafunktion genau solche
universell dualit\"atsstabilen Resonanzen unter nat\"urlichen Lifts in $\Hh$ bzw.\ $\Okt$ sind.
Damit wird ``Satz + Vermutung'' zu einem klaren, logisch transparenten Programm.
\end{abstract}

\section{Motivation und Claim}
Dieses Paper macht \emph{keine} physikalische Behauptung \"uber Strahlungsspektren.
Der Claim ist rein algebraisch-geometrisch:

\begin{quote}
\emph{Wenn ein (liftbarer) komplexer Parameter in einer normierten Divisionsalgebra die Norm-Dualit\"at
$x\mapsto 1-x$ erf\"ullt, dann liegt sein skalarer Anteil zwingend auf $1/2$.}
\end{quote}

Das ist die mathematische Form dessen, was im fr\"uheren Manuskript als ``Hurwitz-Imperativ''
auftauchte, nun jedoch als sauberer Satz mit expliziten Voraussetzungen.

\section{Normierte Divisionsalgebren und komplexe Slices}

\begin{definition}[Normierte Divisionsalgebra]
Eine (reelle) Algebra $\mathbb{A}$ hei\ss t \emph{normierte Divisionsalgebra}, wenn es eine Norm
$\|\cdot\|$ gibt mit \emph{Kompositionsgesetz} $\|xy\|=\|x\|\,\|y\|$ f\"ur alle $x,y\in\mathbb{A}$
und jedes $x\neq 0$ invertierbar ist.
Nach Hurwitz gilt $\mathbb{A}\in\{\R,\C,\Hh,\Okt\}$.
\end{definition}

\begin{definition}[Komplexer Slice]
Sei $\mathbb{A}\in\{\Hh,\Okt\}$ und $u\in\mathrm{Im}(\mathbb{A})$ mit $u^2=-1$.
Dann ist
\[
\C_u := \{\,\sigma + \gamma u : \sigma,\gamma\in\R\,\}\subset\mathbb{A}
\]
ein Unterraum, der als Algebra isomorph zu $\C$ ist (``complex slice'').
\end{definition}

\begin{definition}[Slice-Einbettung]
Die kanonische Einbettung in den Slice sei
\[
\iota_u:\C\to\mathbb{A},\qquad \iota_u(\sigma+\ii\gamma)=\sigma+\gamma u.
\]
\end{definition}

\begin{bemerkung}
Die Norm in $\mathbb{A}$ reduziert sich auf dem Slice auf den euklidischen Betrag:
\[
\|\iota_u(\sigma+\ii\gamma)\|^2 = \sigma^2+\gamma^2 = |\sigma+\ii\gamma|^2.
\]
Dies ist elementar; die inhaltliche Frage liegt nicht in dieser Identit\"at, sondern darin,
\emph{wann} und \emph{f\"ur welche} Parameter eine Dualit\"ats-Normbedingung physikalisch oder
zahlentheoretisch motiviert ist.
\end{bemerkung}

\section{Dualit\"ats-Normfunktional und Fixpunkt}

\begin{definition}[Dualit\"ats-Funktional]
F\"ur $x\in\mathbb{A}$ definieren wir
\[
\Phi(x) := \big|\|x\|^2-\|1-x\|^2\big|.
\]
Auf einem Slice schreiben wir auch $\Phi_u(s):=\Phi(\iota_u(s))$.
\end{definition}

\begin{lemma}[Fixpunkt der Norm-Dualit\"at]\label{lem:fixpoint}
Sei $\mathbb{A}\in\{\R,\C,\Hh,\Okt\}$ und $x=\sigma+\gamma u\in\C_u\subset\mathbb{A}$.
Wenn $\Phi(x)=0$, dann gilt $\sigma=\tfrac12$.
\end{lemma}

\begin{proof}
Auf dem Slice gilt $\|x\|^2=\sigma^2+\gamma^2$ und $\|1-x\|^2=(1-\sigma)^2+\gamma^2$.
Aus $\Phi(x)=0$ folgt $\sigma^2+\gamma^2=(1-\sigma)^2+\gamma^2$, also $2\sigma=1$.
\end{proof}

\begin{satz}[Kritische Achse als Dualit\"ats-Fixpunkt]\label{thm:critical_axis}
In jeder normierten Divisionsalgebra $\mathbb{A}\in\{\R,\C,\Hh,\Okt\}$ liegt jedes Element
$x\in\C_u$ mit $\|x\|=\|1-x\|$ auf der Fixpunktachse $\Re(x)=1/2$.
\end{satz}

\begin{proof}
Unmittelbar aus Lemma~\ref{lem:fixpoint}.
\end{proof}

\section{Universelle Stabilit\"at \"uber Slices}
Das Theorem ist elementar, aber sein Nutzen entsteht, wenn man eine \emph{Universality}-Annahme
formuliert: dass bestimmte ``Resonanzen'' gegen Wahl des Slices stabil sind.

\begin{definition}[Universelle Dualit\"atsstabilit\"at]\label{def:universal}
Sei $\mathbb{A}\in\{\Hh,\Okt\}$ und $\mathcal{U}$ eine Klasse zul\"assiger imagin\"arer Einheiten
$u$ (z.B.\ alle $u$ mit $u^2=-1$ oder eine durch Strukturbedingungen eingeschr\"ankte Teilmenge).
Ein Parameter $s\in\C$ hei\ss t \emph{universell dualit\"atsstabil} (bez\"uglich $\mathcal{U}$), wenn
\[
\Phi(\iota_u(s))=0\quad\text{f\"ur alle }u\in\mathcal{U}.
\]
\end{definition}

\begin{korollar}[Universality $\Rightarrow$ kritische Linie]
Jeder universell dualit\"atsstabile Parameter $s$ erf\"ullt $\Re(s)=1/2$.
\end{korollar}

\section{Vermutung zur Riemannschen Vermutung (klar markiert)}
Wir trennen streng zwischen dem Satzteil (bewiesen) und der Br\"ucke zur Zetafunktion (offen).

\begin{conjecture}[Division-Algebra Universality der nichttrivialen Nullstellen]
Seien $\rho$ nichttriviale Nullstellen der vollendeten Zetafunktion (oder $\xi$-Funktion).
Dann sind diese $\rho$ universell dualit\"atsstabil im Sinn von
Definition~\ref{def:universal} f\"ur eine nat\"urliche Klasse $\mathcal{U}$ von Slices, die durch den
Lift/Analytische Fortsetzung der Riemann-Funktion in $\Hh$ bzw.\ $\Okt$ kanonisch bestimmt ist.
\end{conjecture}

\begin{korollar}[Programm: Satz + Vermutung $\Rightarrow$ RH]
Gilt die obige Vermutung, dann folgt f\"ur jede nichttriviale Nullstelle $\rho$:
$\Re(\rho)=1/2$.
\end{korollar}

\section{Offene Punkte (was zu zeigen w\"are)}
\begin{itemize}
\item Eine kanonische (nichttriviale) Definition des ``Lifts'' der analytisch fortgesetzten Zeta-/
$\xi$-Funktion in $\Hh$ und $\Okt$ (z.B.\ Slice-regular oder Stem-function Framework).
\item Eine begr\"undete Wahl der Slice-Klasse $\mathcal{U}$, die die Zahlentheorie (Funktionalgleichung)
und die Algebra (Automorphismen/Isometrien der Norm) zugleich reflektiert.
\item Der Nachweis, dass die nichttrivialen Nullstellen die Universality-Bedingung erf\"ullen.
\end{itemize}

\section*{Hinweis zum Folgepaper (Paper B)}
Der physikalische Teil (Strahlungsmodulation, Template-Matching, Metrologie) wird bewusst in einem
separaten Manuskript behandelt, um die logische Struktur dieses Papers rein mathematisch zu halten.

\end{document}
